% Volume III: Attention and Publicity
% Part VI. Privacy as Boundary and Capacity
% Chapter 33: The Right to Remain Underdescribed
% Spine: proof-spines/ch033-spine.md

\chapter{The Right to Remain Underdescribed}
\label{ch:ch033}

The chapter introduces \textbf{underdescription} as a political good. Human freedom partly depends on not being perfectly resolved into a predictive object, and a person should not be forced to exhaustively explain every anomaly merely because extensive records make such explanation technically demandable. The relevant right is not a right to deceive; it is a limit on the presumption that whatever can be inferred may therefore be legitimately used.

Placed after model capture, the point is sharper: compulsory completeness turns persons into administratively convenient dossiers. A right to remain underdescribed preserves revisability, experimentation, and context-specific self-presentation by refusing the demand that every trace be made permanently commensurable with every institutional purpose. The same protection becomes concealment when officials or powerful organizations claim a comparable right to remain underdescribed while exercising authority over others.
